\section{Results and Discussion}

Our analysis revealed a substantial difference between conditions. As shown in Figure~\ref{fig:results}, participants in the treatment condition responded approximately 50 milliseconds faster than those in the baseline condition, while maintaining accuracy above 90\% in both groups.

\begin{figure}[h]
\centering
\includegraphics[width=0.7\textwidth]{figures/plot.pdf}
\caption{Mean response time by experimental condition. Error bars represent standard error. The treatment condition shows a clear reduction in response latency compared to baseline.}
\label{fig:results}
\end{figure}

These findings support our hypothesis and align with prior work demonstrating that task-specific training can enhance performance \cite{schneider1977controlled}. The preservation of accuracy suggests that the speed improvement does not come at the cost of decision quality.

\subsection{Implications}

The magnitude of the effect—approximately 20\% reduction in response time—is practically significant for applications where rapid decision-making is critical. Future work should investigate whether these gains persist over longer time periods and transfer to related tasks.
